\ifdefined\gkver\else\def\gkver{student}\fi
\documentclass[\gkver]{gaokaozhenti}
\begin{document}

\chapter{2004年全国综合能力测试}
%% paperType: 综合能力测试物理部分

\begin{enumerate}
\item
%% number: 31
%% paperName: 2004年全国综合能力测试第 31 题 6 分
%% typeId: 1
%% type: 单选题
%% chapter: 原子和原子核
%% point: 原子核的组成
%% method: 无
%% score: 6
%% degree: 850
%% duplicateId: 0
%% body:
若用 $x$ 代表一个中性原子中核外的电子数，$y$ 代表此原子的原子核内的质子数，$z$ 代表此原子的原子核内的中子数，则对 \ce{^{234}_{90}Th} 的原子来说\xzanswer{B}

\fourchoices[answer=B]
{$x=90$，$y=90$，$z=234$}
{$x=90$，$y=90$，$z=144$}
{$x=144$，$y=144$，$z=90$}
{$x=234$，$y=234$，$z=324$}

%% answer: B
%% memo:
\memoanswer{
本题原卷及材料中未提供详解，待补充。
}

\item
%% number: 32
%% paperName: 2004年全国综合能力测试第 32 题 6 分
%% typeId: 1
%% type: 单选题
%% chapter: 运动和力
%% point: 牛顿第二定律
%% method: 无
%% score: 6
%% degree: 750
%% duplicateId: 0
%% body:
三个完全相同的物块 1、2、3 放在水平桌面上，它们与桌面间的动摩擦因数都相同。现用大小相同的外力 $F$ 沿图示方向分别作用在 1 和 2 上，用 $\frac{1}{2}F$ 的外力沿水平方向作用在 3 上，使三者都做加速运动。令 $a_{1}$、$a_{2}$、$a_{3}$ 分别代表物块 1、2、3 的加速度，则\xzanswer{C}

\onepicture[w=6.59cm]{figs/32.png}

\fourchoices[answer=C]
{$a_{1}=a_{2}=a_{3}$}
{$a_{1}=a_{2}$，$a_{2}>a_{3}$}
{$a_{1}>a_{2}$，$a_{2}<a_{3}$}
{$a_{1}>a_{2}$，$a_{2}>a_{3}$}

%% answer: C
%% memo:
\memoanswer{
本题原卷及材料中未提供详解，待补充。
}

\item
%% number: 33
%% paperName: 2004年全国综合能力测试第 33 题 6 分
%% typeId: 1
%% type: 单选题
%% chapter: 光学
%% point: 光子说
%% method: 无
%% score: 6
%% degree: 750
%% duplicateId: 0
%% body:
下列说法中正确的是\xzanswer{C}

\fourchoices[answer=C]
{在真空中红光的波长比紫光的小}
{玻璃对红光的折射率比对紫光的大}
{在玻璃中红光的传播速度比紫光的大}
{红光光子的能量比紫光光子的能量大}

%% answer: C
%% memo:
\memoanswer{
本题原卷及材料中未提供详解，待补充。
}

\item
%% number: 34
%% paperName: 2004年全国综合能力测试第 34 题 6 分
%% typeId: 1
%% type: 单选题
%% chapter: 功和能
%% point: 动能定理求路程
%% method: 无
%% score: 6
%% degree: 650
%% duplicateId: 0
%% body:
如图所示，ABCD 是一个盆式容器，盆内侧壁与盆底 BC 的连接处都是一段与 BC 相切的圆弧，B、C 为水平的，其距离 $d=0.50\Um$，盆边缘的高度为 $h=0.30\Um$。在 A 处放一个质量为 $m$ 的小物块并让其从静止出发下滑。已知盆内侧壁是光滑的，而盆底 BC 面与小物块间的动摩擦因数为 $\mu=0.10$。小物块在盆内来回滑动，最后停下来，则停的地点到 B 的距离为\xzanswer{D}

\onepicture[w=5.74cm]{figs/34.png}

\fourchoices[answer=D]
{$0.50\Um$}
{$0.25\Um$}
{$0.10\Um$}
{$0$}

%% answer: D
%% memo:
\memoanswer{
本题原卷及材料中未提供详解，待补充。
}

\item
%% number: 35
%% paperName: 2004年全国综合能力测试第 35 题 6 分
%% typeId: 1
%% type: 单选题
%% chapter: 电磁感应
%% point: 电磁感应受力分析
%% method: 无
%% score: 6
%% degree: 750
%% duplicateId: 0
%% body:
如图所示，ab 和 cd 是位于水平面内的平行金属轨道，其电阻可忽略不计。ac 之间连接一阻值为 $R$ 的电阻。ef 为一垂直于 ab 和 cd 的金属杆，它与 ab 和 cd 接触良好并可沿轨道方向无摩擦地滑动。ef 长为 $l$，电阻可忽略。整个装置处在匀强磁场中，磁场方向垂直于图中纸面向里，磁感应强度为 $B$，当施外力使杆 ef 以速度 $v$ 向右匀速运动时，杆 ef 所受的安培力为\xzanswer{A}

\onepicture[w=5.58cm]{figs/35a.png}

\fourchoices[answer=A]
{$\frac{vB^2l^2}{R}$}
{$\frac{vBl}{R}$}
{$\frac{vB^2l}{R}$}
{$\frac{vBl^2}{R}$}

%% answer: A
%% memo:
\memoanswer{
本题原卷及材料中未提供详解，待补充。
}

\item
%% number: 35
%% paperName: 2004年全国综合能力测试第 35 题 6 分
%% typeId: 1
%% type: 单选题
%% chapter: 磁场
%% point: 带电粒子在磁场中的运动
%% method: 无
%% score: 6
%% degree: 950
%% duplicateId: 0
%% body:
（新课程）如图为示波管中偏转电极的示意图，相距为 $d$、长度为 $l$ 的平行板 A、B 加上电压后，可在 A、B 之间的空间中（设为真空）产生电场（设为匀强电场）。在 AB 左端距 A、B 等距离处的 O 点，有一电量为 $+q$、质量为 $m$ 的粒子以初速 $v_{0}$ 沿水平方向（与 A、B 板平行）射入（如图）。不计重力，要使此粒子能从 C 处射出，则 A、B 间的电压应为\xzanswer{A}

\onepicture[w=5.95cm]{figs/35b.png}

\fourchoices[answer=A]
{$\frac{d^2mv_{0}^{2}}{ql^2}$}
{$\frac{l^2mv_{0}^{2}}{qd^2}$}
{$\frac{lmv_{0}}{qd}$}
{$\frac{qmv_{0}}{dl}$}

%% answer: A
%% memo:
\memoanswer{
本题原卷及材料中未提供详解，待补充。
}

\item
%% number: 36
%% paperName: 2004年全国综合能力测试第 36 题 6 分
%% typeId: 4
%% type: 实验题
%% chapter: 电路
%% point: 测电动势与内阻
%% method: 无
%% score: 6
%% degree: 750
%% duplicateId: 0
%% body:
可用理想电压表 V、理想电流表 A、变阻器 $R$ 以及电键 K 和导线等器材来测量某一电源 $E$ 的电动势和内阻。下面给出了四个电路图，图中 $+$、$-$ 代表电源的正、负极和电表的正负接线柱。正确的电路图是\xzanswer{B}

\fourchoices[answer=B, ispicture=true, h=3.0cm]
{figs/36a.png}
{figs/36b.png}
{figs/36c.png}
{figs/36d.png}

%% answer: B
\jdanswer{
B
}
%% memo:
\memoanswer{
本题原卷及材料中未提供详解，待补充。
}

\item
%% number: 43
%% paperName: 2004年全国综合能力测试第 43 题 6 分
%% typeId: 3
%% type: 填空题
%% chapter: 运动和力
%% point: 超重失重
%% method: 无
%% score: 6
%% degree: 550
%% duplicateId: 0
%% body:
飞船降落过程中，在离地面高度为 $h$ 处速度为 $v_{0}$，此时开动反冲火箭，使船开始做减速运动，最后落地时的速度减为 $v$。若把这一过程当作为匀减速运动来计算，则其加速度的大小等于\tkanswer[2.5cm]{$\frac{v_{0}^{2}-v^{2}}{2h}$}。

已知地球表面处的重力加速度为 $g$，航天员的质量为 $m$，在这过程中航天员对坐椅的压力等于\tkanswer[3.0cm]{$mg+\frac{m(v_{0}^{2}-v^{2})}{2h}$}。
%% answer: $\frac{v_{0}^{2}-v^{2}}{2h}$，$mg+\frac{m(v_{0}^{2}-v^{2})}{2h}$
%% memo:
\memoanswer{
本题原卷及材料中未提供详解，待补充。
}

\end{enumerate}

\end{document}
